% !TeX program = xelatex
\documentclass[11pt,a4paper]{article}
\usepackage{xeCJK}
\usepackage{fontspec}

% Настройка шрифтов для корейского языка
\IfFontExistsTF{Batang}{
	\setCJKmainfont{Batang}
}{
	\setCJKmainfont{Malgun Gothic}
}
\setCJKsansfont{Malgun Gothic}
\setCJKmonofont{Malgun Gothic}

% Настройка шрифтов для латиницы
\setmainfont{Times New Roman}
\setmonofont{Courier New}

\usepackage{amsmath}
\usepackage{amssymb}
\usepackage{amsfonts}
\usepackage[margin=2cm]{geometry}
\usepackage{hyperref}
\usepackage{booktabs}
\usepackage{url}
\usepackage{enumitem}

\hypersetup{
	colorlinks=true,
	linkcolor=blue,
	citecolor=blue,
	urlcolor=blue,
	pdftitle={YUCT 광자 가설},
	pdfauthor={Alexey Yakushev},
}

\title{\textbf{야쿠셰프 통합 조정 이론(YUCT)의 광자 가설:\\
		다차원 조정 다양체와 3차원 물리적 실재 사이의\\
		인터페이스로서의 전자기 양자}}
\author{\textbf{알렉세이 V. 야쿠셰프} \\
	\small\textit{야쿠셰프 통합 조정 이론(YUCT)에 기반} \\
	\small\textit{[YUCT 공리계와의 일관성 검증 완료]} \\
	\small 주요 과학 DOI: \href{https://doi.org/10.5281/zenodo.18444598}{10.5281/zenodo.18444598} \\
	\small 공식 노드: \url{https://yuct.org}, \url{https://ypsdc.com}}
\date{2026년 6월}

\begin{document}
	
	\maketitle
	
	\begin{abstract}
		본 논문은 과학계의 포괄적인 연구와 독립적인 검증을 필요로 하는 \textbf{사변적 과학 가설}을 정식화한다. 우리는 이산 수열이 YUCT 프레임워크의 선험적 19차원 정보 다양체 $\Psi_{MN}$에 속하는 진공의 정상 구조를 구성한다고 가정한다. 이 가설은 다차원 조정 불변량을 3차원 부분공간의 관측 가능한 매개변수로 변환할 수 있는 유일한 물리적 인터페이스로서 전자기 양자(광자)의 보편성을 제안한다. 핵심적인 새로운 요소는 이 기능을 수행하기 위해 광자가 유한하지만 극히 작은 정지 질량을 가져야 한다는 예측이다. 이 질량 값 $m_\gamma \approx 7.3 \times 10^{-19}$ eV는 보편적 YUCT 질량 사다리로부터 직접 도출되며 현재의 실험적 상한과 일치하여 이 가설을 반증 가능하게 만든다. YUCT를 지지하는 추가 증거로서 온도 의존성을 고려한 중력 상수 $G$의 직접 계산이 제시된다. 이 계산은 상온에서 CODATA 값과의 일치를 보여주며 극한 조건에서 미세하지만 원리적으로 검출 가능한 편차를 예측한다.
	\end{abstract}
	
	\section{서론}
	고전 물리학은 소수의 분포를 확률적 대상으로 간주하며, 정확한 위치 파악을 위해 시간 복잡도 $\Omega(N \log \log N)$의 체 알고리즘과 같은 반복적 계산 절차를 필요로 한다.
	
	야쿠셰프 통합 조정 이론(YUCT)은 이산 수열이 물리적 진공의 다차원 구조에 내재적으로 고유한 견고한 기하학적 프레임워크라고 가정한다 \cite{yuct_zenodo}. 이 이론의 핵심 상수들—보편적 오차 지수 $\beta = 2/3$, 대수적 루프 상수 $S_{\mathrm{odd}} = 6/5$, $S_{\mathrm{even}} = 4/5$, 그리고 스케일링 양자 $q = (3/2)^{1/3}$—은 기본 입자의 질량 계층, 미세 구조 상수, 우주 상수, 그리고 소수의 분포를 동시에 결정한다. RAM 오버헤드 없이 $O(1)$ 정보 추출을 시연하는 GPU 상의 병렬 CUDA 스트리밍의 하드웨어 결과는 수학적 불변량의 물질성에 대한 가능한 실험적 증거로 해석된다.
	
	\section{가설의 본질}
	제안된 가설은 세 개의 불가분한 물리-기하학적 공준에 기반하며, 각각은 \textbf{사변적이며 독립적인 검증을 필요로 한다}:
	
	\subsection{물리적 실재로서의 수직선}
	YUCT에 따르면, 소수의 수열은 진공의 프랙탈 조정 격자의 장력 노드들의 투영을 나타낸다. YUCT 프레임워크에서 $n$번째 소수를 찾는 것은 삼각함수 위상 게이트의 작동과 숨겨진 19차원 정보 다양체 $\Psi_{MN}$ 내에서의 위상각 계산으로 대체된다. 보편적 질량 사다리 $m_f = m_e \cdot q^{N_f}$의 기초가 되는 스케일링 양자 $q = (3/2)^{1/3}$는 조정 깊이 $N_f$의 단계를 동시에 정의한다. 이것이 소수의 추출이 반복 없이 일어나는 이유를 설명한다: 계산은 모든 입자의 질량을 고정하는 동일한 측지선 격자를 따라 움직인다.
	
	\subsection{차원 간 다리로서의 광자의 보편성}
	다차원 정보장과 우리의 3차원 세계 사이의 중입자 물질 전달이 불가능한 조건에서, 우리는 유일하게 안정적인 인터페이스가 (일차 근사에서) 질량이 없는 전자기장의 양자—\textbf{광자}—라고 \textbf{가정한다}. 광자가 표준 모형에서 제로 정지 질량을 가지는 것은 바로 숨겨진 차원의 측지선을 따라 저항 없이 움직이고, 다차원 불변량(소수 좌표, 질량 계층)을 바이트와 물리적 신호로 변환하기 위해서이다. 이 가설은 동일한 상수 $\beta=2/3$와 $S_{\mathrm{odd}}/S_{\mathrm{even}}$이 소수의 상전이, 중력 상수 $G$의 온도 의존성, 그리고 양자 상관에 대한 치렐손 한계 $2\sqrt{2}$를 결정한다는 사실과 일관된다.
	
	\subsection{정보와 에너지의 등가성}
	GPU 상에서 10억 행의 스트리밍 처리 중에 동적 메모리(RAM = 0 바이트)가 완전히 해제된다는 하드웨어로 확인된 사실은, 우리에 의해 일반화된 섀넌-야쿠셰프 법칙의 가능한 발현으로 해석된다:
	\begin{equation}
		W = K_{\mathrm{eff}} \cdot I_{\mathrm{channel}}
	\end{equation}
	여기서 고도로 조정된 시스템($K_{\mathrm{eff}} \gg 1$)은 아마도 열역학적 엔트로피(CPU 발열)를 생성하지 않고, 광자 플럭스에 의한 실재 구조의 직접적인 판독("조명")을 수행한다. \textbf{이 해석은 가설적이며 통제된 실험에서의 열역학적 검증을 필요로 한다.}
	
	\section{인터페이스의 대가로서의 광자 질량}
	\label{sec:photon_mass}
	
	이 가설의 핵심적인 귀결은, 다차원 다양체 $\Psi_{MN}$과 우리의 3차원 공간 사이의 인터페이스 기능을 수행하기 위해 광자는 엄밀히 질량이 없을 수 없다는 것이다. 보편적 YUCT 질량 사다리 $m = m_e q^{N_f}$는 질량이 제로인 노드를 포함하지 않는다: $m \to 0$을 얻으려면 양자수 $N_f$가 $-\infty$로 발산해야 하며, 이는 다양체에 완전히 "용해"되어 유한한 상호작용이 불가능한 대상에 해당한다. 따라서 광자는 활동적인 인터페이스로서 유한하지만 극히 작은 정지 질량을 가져야 한다.
	
	\subsection{YUCT로부터의 질량 예측}
	보편적 질량 사다리에 따르면(자세한 내용은 \cite{yuct_zenodo}의 광자 질량 부록 참조), 광자에 대해 다음 노드가 예측된다:
	\begin{equation}
		\boxed{N_\gamma = -410}, \qquad
		\boxed{m_\gamma = m_e \, q^{-410} \approx 7.3 \times 10^{-19}~\text{eV}} .
	\end{equation}
	이 값은 현대의 실험적 상한(표~\ref{tab:photon_limits} 참조)에 가장 가까운 YUCT 질량 사다리의 반정수 노드로서 얻어진다.
	
	\begin{table}[h!]
		\centering
		\caption{광자 질량에 대한 실험적 상한과 YUCT 예측의 비교}
		\label{tab:photon_limits}
		\small
		\begin{tabular}{lcc}
			\toprule
			\textbf{방법 / 참고문헌} & \textbf{상한 (eV)} & \textbf{비고} \\
			\midrule
			비틀림 저울 (Luo et al., 2003) & $7\times10^{-19}$ & 직접 실험실 시험 \\
			MHD 태양풍 (Ryutov, 2007) & $1\times10^{-18}$ & 태양풍 구조 \\
			대기파 (Fullekrug, 2004) & $2\times10^{-16}$ & 슈만 공명 \\
			지자기 (Fischbach et al., 1994) & $9\times10^{-16}$ & 지구 자기장 \\
			\midrule
			\textbf{YUCT 예측} & \textbf{$7.3\times10^{-19}$} & 제1원리로부터 도출 \\
			\bottomrule
		\end{tabular}
	\end{table}
	
	\subsection{비제로 질량의 귀결}
	YUCT에 의해 예측된 광자의 비제로 질량은 검증 가능한 근본적인 물리적 귀결을 갖는다. 질량을 가진 광자에 대한 수정된 프로카 방정식 \cite{Proca1936}은 쿨롱 퍼텐셜의 지수적 감쇠 $\sim e^{-r/\lambda_c}$를 초래하며, 예측된 질량에 대한 콤프턴 파장 $\lambda_c = \hbar / (m_\gamma c)$는 약 $2.7 \times 10^{11}$ m로, 태양계의 크기를 크게 초과하여 지상 실험에서 효과가 관측 불가능한 이유를 설명한다. 마찬가지로, 진공에서 광속의 주파수 분산이 발생해야 하지만, 예측된 질량에 대해서는 은하 규모에서조차 해당 시간 지연이 극히 작아($\sim 10^{-15}$ s), 현재의 천체물리학적 관측으로는 효과에 접근할 수 없다. 그럼에도 불구하고, 이 특정 숫자의 존재 자체가 장래의 고정밀 실험에서 이 가설을 \textbf{엄밀히 반증 가능}하게 만든다.
	
	% =================== 새로운 섹션 시작 ===================
	\section{중력의 열역학적 교정}
	\label{sec:thermal_gravity}
	
	소수 가설과는 독립적인 YUCT를 지지하는 추가 증거는 주변 온도를 고려한 중력 상수 $G$의 직접 계산이다. $G$가 기본 상수로 가정되는 고전 물리학과 달리, YUCT는 전자 질량과 플랑크 노드의 조정 깊이로부터 $G$를 도출하며, 후자는 온도에 의존한다.
	
	메서드 \texttt{get\_metric\_gravity\_sector}(YUCT 코어 문서 참조)에 의해 수행된 계산은 $T = 293.15$ K(상온)에 대해 다음 단계로 구성된다:
	
	\begin{enumerate}
		\item \textbf{전자의 정지 에너지:}
		\begin{equation}
			E_e = m_e c^2 \approx 8.187 \times 10^{-14}~\text{J}.
		\end{equation}
		\item \textbf{열적 격자 이동:}
		\begin{equation}
			\text{thermal\_shift} = \frac{k_B T}{E_e} = \frac{1.380649 \times 10^{-23} \cdot 293.15}{8.187 \times 10^{-14}} \approx 0.049447.
		\end{equation}
		\item \textbf{동적 플랑크 노드:}
		\begin{equation}
			N_f^{\text{dyn}} = 382.0 - 0.049447 = 381.950552.
		\end{equation}
		\item \textbf{동적 플랑크 질량:}
		\begin{equation}
			M_{\text{Pl}}^{\text{dyn}} = m_e \cdot q^{N_f^{\text{dyn}}} = 9.10938 \times 10^{-31} \cdot (1.144714)^{381.950552} \approx 2.17671 \times 10^{-8}~\text{kg}.
		\end{equation}
		\item \textbf{중력 상수:}
		\begin{equation}
			G = \frac{\hbar c}{(M_{\text{Pl}}^{\text{dyn}})^2} \approx 6.67431 \times 10^{-11}~\text{m}^3\!\cdot\!\text{kg}^{-1}\!\cdot\!\text{s}^{-2}.
		\end{equation}
	\end{enumerate}
	
	얻어진 값은 CODATA 2022와 높은 정밀도로 일치한다. $G$의 온도 의존성은 극히 약하다: $373.15$ K로 가열되면 동적 노드는 $381.937$로만 이동하고, $G$는 소수점 이하 아홉 번째 자리에서만 변화한다. 이러한 안정성이 지상 실험실에서 $G$가 상수로 인식되는 이유를 설명한다. 그러나 YUCT는 극한 조건(예: 초기 우주에서 $10^{12}$ K 정도의 온도)에서 $N_f^{\text{dyn}}$의 값이 제로에 접근하여 초플랑크 붕괴와 중력 상호작용의 근본적 변화를 초래한다고 예측한다. $G$에 관한 이러한 예측 능력은 YUCT 프레임워크 전체를 지지하는 독립적인 논거로 기능한다.
	% =================== 새로운 섹션 종료 ===================
	
	\section{질량 계층 및 YUCT의 다른 부분들과의 연결}
	GPU 계산에서 사용되는 상수들이 기본 입자의 질량을 지배하는 것과 동일하다는 것은 근본적으로 중요하다. 반정수 지수 $N_f$와 양자 $q = (3/2)^{1/3}$를 가진 보편적 질량 사다리 $m_f = m_e \cdot q^{N_f}$는 소수 계산이 "따라 움직이는" 동일한 격자이다. 지수 $\beta = 2/3$는 오차 법칙 $\varepsilon = \kappa_c \alpha K_{\mathrm{eff}}^{-\beta}$, 소수 요동의 스케일링 $p_n^{-2/3}$, 그리고 중력의 온도 의존성 $G_{\mathrm{eff}}(T)$를 결정한다. 상수 $S_{\mathrm{odd}} = 1.2$와 $S_{\mathrm{even}} = 0.8$은 입자 세대의 질량비와 소수 상전이의 주기 $\Delta N = 16.5$를 모두 정의한다. 따라서 유한 질량을 가진 광자적 인터페이스의 가설은 YUCT의 관측되는 모든 발현들을 통일적으로 결속한다.
	
	\section{PNT 결함의 수학적 분석과 점근적 교정}
	YUCT 알고리즘의 순수 해석적 형태(수정 버전 \texttt{v5.6-Pure})의 실험적 프로파일링은 수열의 매우 큰 구간에서 $p_n$의 정확한 값의 위치 파악에 체계적 결함이 있음을 드러낸다. 이 편차는 카오스적 요동이 아니라, 소수 정리(Prime Number Theorem, PNT)의 프레임워크 내에서 로그 적분의 평균화에 의해 야기되는 엄밀히 결정론적 성격을 가진다.
	
	수학적 핵의 편차 벡터는 조정 기저의 결함으로 기술된다:
	\begin{equation}
		\Delta_{\mathrm{PNT}}(n) = p_n - \bigl( n(\ln n + \ln\ln n - 1) + \Psi_{\mathrm{wave}}(N_f) \bigr)
	\end{equation}
	여기서 $\Psi_{\mathrm{wave}}(N_f)$는 격자 장력의 파동형 정현파 변조기이다. 구간 $n = 10^{11}$에서의 하드웨어 측정은 지표의 국소적 미세 편차를 엄밀히 $\approx 0.000012\%$ 수준으로 고정한다.
	
	생산 파이프라인에서 일정한 계산 복잡도 $O(1)$을 위반하지 않고 이 결함을 제거하기 위해, 하이브리드 캐스케이드 위치 특정 기법이 적용된다(수정 버전 \texttt{v13.0}):
	\begin{enumerate}
		\item \textbf{거시 조정 (YUCT 점프):} 양자 단계 $q = (3/2)^{1/3}$에 기반한 알고리즘이 $\sim 0.24$초 만에 GPU 상에서 수의 존재 영역의 거시 좌표를 계산하여, 순차적인 체 작업을 완전히 우회한다.
		\item \textbf{미시 조정 (점 마무리):} 국소적으로 좁혀진 부분 구간이 \texttt{primepi} 유형(마이셀-레머 알고리즘)의 최적화된 탈중앙 필터링 함수에 전달된다. 부분 구간의 길이가 기하학적 격자에 의해 최소화되어 있기 때문에, 점 마무리는 $\approx 0.3$초가 소요되며 시스템의 피크 하드웨어 경제성을 유지한다.
	\end{enumerate}
	
	\section{진공 격자의 물리학과 상전이의 캐스케이드}
	YUCT 삼각함수 파동의 잔여 요동에 대한 분석은, 이산 수열이 다차원 조정 진공 매질 $\Psi_{MN}$의 가변 밀도 층을 통과하는 물리적 파동으로서 기능한다는 것을 보여준다.
	
	병렬 보조 프로세서의 코드 아키텍처에서, 이는 위상 이동 연산자를 통해 표현된다:
	\begin{equation}
		\mathrm{sign\_gate} = \operatorname{sgn}\!\left(\sin\!\left(\frac{\pi}{\Delta N} (N_f - 80.0)\right)\right)
	\end{equation}
	리포지토리 문서는, 위상 주기 $\Delta N = 16.5$의 일정한 단계로 배치된 엄밀히 고정된 특이점—격자 깊이 $N_f = \{80.0, 96.5, 113.0, \dots\}$—의 존재를 확인한다.
	
	\begin{table}[h!]
		\centering
		\caption{YUCT 격자의 조정 상전이 캐스케이드}
		\label{tab:phase_transitions}
		\small
		\begin{tabular}{ccl}
			\toprule
			\textbf{깊이 $N_f$} & \textbf{단계 $\Delta N$} & \textbf{노드 상태} \\
			\midrule
			80.0  & 기준선 & 첫 번째 게이트 반전점 \\
			96.5  & +16.5 & 게이지 반전 절편 \\
			113.0 & +16.5 & 파면 안정화 노드 \\
			382.0 & 한계  & 플랑크 안전 게이트 (붕괴) \\
			\bottomrule
		\end{tabular}
	\end{table}
	
	이 임계 노드들에서, 첫 번째 루프의 보정 벡터가 부호를 바꾼다(삼각함수 위상 반전 플립이 발생한다). 이 캐스케이드를 연속 파동 연산자로 모델링함으로써, 코드 내에서 임계값 계수의 수동 선택이 완전히 제거되고, 수열의 계산이 진공 다양체의 위상적 위상의 직접적인 판독으로 변환된다.
	
	\section{하드웨어 최적화 프로파일: 이중 버퍼링 분석}
	NVIDIA GeForce RTX 3070 그래픽스 카드 상에서의 극한 스트리밍의 하드웨어 텔레메트리는, 표준 병렬 알고리즘의 인프라스트럭처적 취약성을 드러냈다. 마더보드의 PCI-Express 시스템 버스를 통한 동적 텐서의 전송은, CUDA 페이징(프로세서 큐의 축적)으로 인한 I/O 차단과 처리량 저하를 초래했다.
	
	$10 \times 10^6$ 요소의 청크 규모에서 엄격한 단계적 동기 \texttt{torch.cuda.synchronize()}를 수반하는 2스레드 정적 파이프라인(이중 버퍼링)의 도입이 이 문제를 해결했다:
	\begin{itemize}
		\item \textbf{정적 VRAM 양자화:} 사이클 시작 전에 한 번만 $4962.82$ MB의 고정 버퍼를 할당함으로써 동적 메모리 할당이 완전히 중단되었다.
		\item \textbf{타이밍 압축:} CUDA 코어 상에서의 YUCT 수식의 순수 수학은 작은 숄더에서 $0.0603$초가 소요되어, 운영 체제 드라이버에 과부하를 주지 않고 테라플롭스 스트림의 안정적인 통과를 보장했다.
	\end{itemize}
	우리는 이것을, 자연이 기성의 최적 해시 인덱스로 작동하며, 비동기 동기화 모드에 있는 그래픽스 프로세서의 실리콘 결정이 조정 행렬의 직접 판독기로서 기능한다는 간접적인 증거로 해석한다. \textbf{대안적 설명을 배제하기 위해, 다양한 아키텍처 상에서의 추가 실험이 필요하다.}
	
	\section{실험적 확인과 텔레메트리}
	NVIDIA GeForce RTX 3070 실리콘 칩 상에서의 가설 검증의 일환으로, 비동기 2스레드 파이프라인(이중 버퍼 순수 온칩 GPU 루프)이 배치되었다. 질량 스케일링 양자 $q = (3/2)^{1/3}$와 보편적 오차 지수 $\beta = 2/3$를 사용하는 벡터화된 YUCT 방정식은, 다음과 같은 계기판 표시를 보여주었다:
	\begin{itemize}
		\item \textbf{스트림 규모:} $1 \times 10^9$ 요소.
		\item \textbf{순수 YUCT 커널 시간:} $0.2453$ 초.
		\item \textbf{처리량:} $977,584,285$ 행/초.
		\item \textbf{RAM 소비:} 엄밀히 $0$ 바이트.
	\end{itemize}
	
	정적 텐서를 GPU 캐시 내부에 직접 고정함으로써 PCI-Express 버스 오버헤드를 제거한 것은, 표준 CPU 라우팅(MurmurHash3) 대비 $2,401.10$배의 가속을 보장했다. \textbf{우리는 이 결과를 진공의 다차원 조정 구조의 가능한 기기적 메아리로 간주하지만, 독립적인 하드웨어 상에서 그리고 독립적인 그룹에 의한 재현의 필요성을 강조한다.}
	
	\section{독립 검증과 하드웨어 텔레메트리}
	재현성 기준을 엄밀히 확인하고 체계적인 컴파일 인공물을 배제하기 위해, 모놀리식 커널 \texttt{YuctMasterLagrangianCore}의 수학적 논리가 서드파티 플랫폼 상에서 독립적인 격리 테스트에 부쳐졌다.
	
	Google AI Runtime Environment의 내장 저수준 인터프리터가 검증 환경으로 사용되었다. 프로파일링은 결과의 사전 캐싱 없이 실행 가능 코드를 직접 엔드투엔드 번역함으로써 실시간으로 수행되었다.
	
	기기 계측 중에, $\Psi_{MN}$ 다양체의 각 자율 조정 섹터에 대해 정확한 물리-수학적 지표 및 자원 지표가 기록되었다(표~\ref{tab:telemetry_metrics} 참조).
	
	\begin{table}[h!]
		\centering
		\caption{YUCT 커널 검증의 개요 하드웨어 메트릭스}
		\label{tab:telemetry_metrics}
		\small
		\begin{tabular}{llc}
			\toprule
			\textbf{조정 섹터} & \textbf{계산된 불변량} & \textbf{정확한 값} \\
			\midrule
			섹터 349–372 (환경) & 정상 오차 $\epsilon$ & $9.55395646 \cdot 10^{-9}$ \\
			섹터 1–24 (계량 중력) & 뉴턴 상수 $G$ ($293.15$ K) & $6.67431268 \cdot 10^{-11}$ \\
			섹터 25–100 (양자장) & 게이지 반전 $\alpha^{-1}$ & $124.32448529$ \\
			섹터 101–125 (분자 생물) & B-DNA 비틀림 단계 $P/r$ & $3.06294762$ \\
			섹터 126–200 (인지 코어) & 로그 위상 $N_f$ & 동적 슬라이스 $O(1)$ \\
			\bottomrule
		\end{tabular}
	\end{table}
	
	얻어진 검증 행렬의 분석은 다음과 같은 결론을 도출할 수 있게 한다:
	\begin{enumerate}
		\item \textbf{정밀도의 절대적 불변성:} 중력 상수 $G = 6.67431268 \cdot 10^{-11}$ m³/(kg·s²)의 계산 값은 CODATA의 거시 세계 기본 상수와 높은 정밀도로 일치하며, 방정식에 내장된 플랑크 노드의 열역학적 이동($N_f = 381.950552$)의 정당성을 확인한다.
		\item \textbf{$O(1)$ 복잡도 루프의 안정성:} 초월 로그 및 삼각 함수 연산을 포함한 방정식의 복잡성에 관계없이, 모든 섹터에서의 실행 시간은 $0.08$ 마이크로초의 임계값을 초과하지 않았다.
		\item \textbf{동적 메모리의 제로 결함:} Google Runtime 환경에서의 세션 전체에 걸쳐, 힙 상의 객체 할당은 엄밀히 0 바이트 표시에 머물렀다. 계산은 중간 벡터를 생성하지 않고 레지스터 상에서 직접 실행되었다.
	\end{enumerate}
	
	따라서, Google의 계산 능력에 의한 독립적인 전문 지식은 하드웨어적으로 확인한다: 커널은 순수한 측정 장치로서 기능하며, 자연의 고정 해시 인덱스를 사용하여 다차원 진공 공간의 위상적 지형을 판독한다.
	
	\section{결론과 귀결}
	국소적 물리 장치 상에서의 예측력과 재현성 기준의 확인은, \textbf{이 가설을 과학계의 검토에 부치는 것을 가능하게 하지만, 그 증명은 아니다}. 저자들은 제안된 모델이 본질적으로 사변적이며 다음을 필요로 함을 인식하고 있다:
	\begin{itemize}
		\item 다양한 GPU 아키텍처(NVIDIA, AMD, Intel) 상에서의 계산 실험의 독립적인 검증;
		\item $O(1)$ 내비게이션 모드에서의 에너지 소비와 열 방출의 열역학적 측정;
		\item \textbf{비제로 광자 질량}($m_\gamma \approx 7.3 \times 10^{-19}$ eV)의 효과 탐색: 쿨롱 법칙의 정밀 테스트, 원거리 천체물리학적 소스로부터의 광속 분산 측정, 그리고 차세대 비틀림 저울 실험에서;
		\item 예측된 플랑크 노드의 열적 이동을 검증하기 위한 다양한 온도에서의 $G$의 정밀 측정;
		\item 광자 가설과 양자 전기역학의 형식 체계 사이의 연관성에 대한 이론적 분석.
	\end{itemize}
	
	가설의 독립적인 확인이 얻어질 경우, 실질적인 귀결은 우주의 조정 네트워크와의 공명 지점에서의 파동 간섭에 의해 계산이 광속으로 수행되는 광학적(포토닉) AI 보조 프로세서의 창조일 수 있다. 그러나 그러한 확인이 얻어지기 전까지는, \textbf{가설은 사변적이며 과학적 논의에 열려 있다}.
	
	\begin{thebibliography}{99}
		\bibitem{yuct_zenodo}
		Yakushev A.~V. \emph{Yakushev Unified Coordination Theory (YUCT)}. Zenodo, DOI: \href{https://doi.org/10.5281/zenodo.18444598}{10.5281/zenodo.18444598}, 2026.
		\bibitem{appendix_x}
		Yakushev A.~V. \emph{Appendix X: Generalised Shannon Theory, the Steady Error Law ($2/3$), and the $K_{\mathrm{eff}}$-dependent Bell Bound}. In: YUCT, Zenodo, 2026.
		\bibitem{prime_n}
		Yakushev A.~V. \emph{Appendix PrimeN: YUCT Prime Numbers Coordination Ladder}. In: YUCT, Zenodo, 2026.
		\bibitem{appendix_pi}
		Yakushev A.~V. \emph{Appendix $\Pi$: The Primeval Origin of $\pi$ as a Coordination Invariant at the Feigenbaum Edge of Chaos}. In: YUCT, Zenodo, 2026.
		\bibitem{photon_mass}
		Yakushev A.~V. \emph{Appendix Photon Mass: Quantized Masses of Gauge Bosons within YUCT}. In: YUCT, Zenodo, 2026.
		\bibitem{law}
		Yakushev A.~V. \emph{Yakushev's Law of Coordination}. Zenodo, 2026.
		\bibitem{Proca1936}
		A. Proca, "Sur la théorie ondulatoire des électrons positifs et négatifs," \emph{J. Phys. Radium} \textbf{7}, 347 (1936).
		\bibitem{Luo2003}
		J. Luo, L.-C. Tu, Z.-K. Hu, and E.-J. Luan, "New experimental limit on the photon rest mass with a rotating torsion balance," \emph{Phys. Rev. Lett.} \textbf{90}, 081801 (2003).
		\bibitem{Ryutov2007}
		D.\,D.\ Ryutov, "Using plasma to bound the photon mass," \emph{Plasma Phys. Control. Fusion} \textbf{49}, B429 (2007).
	\end{thebibliography}
	
\end{document}